\section{Solving the model}\label{sec:q:sol}

\subsection{Terminal conditions}\label{sec:q:sol:terminal}

The six terminal conditions replace the missing period-$T{+}1$ information. All six take one form. If from date $T$ onward the dividend of a costate and the costate itself grow at a common factor $\Gamma$ while the survival factor is constant, its recursion $\left(1+r\right)m_T=\mathrm{div}_{T+1}+s\,m_{T+1}$ closes at
\begin{align}
\boxed{\;m_T \;=\; \frac{\Gamma\,\mathrm{div}_T}{1+r^{\infty}-s_T\,\Gamma}\;},
\qquad
1+r^{\infty} \;=\; \frac{\Gamma^{\eta}}{\beta},
\label{eq:q:terminal}
\end{align}
with $\mathrm{div}$ and $s$ read from Table~\ref{tab:q:costates} at date $T$ and $r^{\infty}$ the balanced-growth interest rate implied by the Euler equation. Three properties make~\eqref{eq:q:terminal} the right closure to use rather than the more familiar ``impose a steady state at $T$''.

\emph{It nests the classification and its benchmark.} Setting $\Gamma=1$ gives the rest-point formulae of the stationary benchmark of \theory{Section}{The floorless benchmark: stationary technology}, including $p^{\mathcal W}=\mu h/\left(\rho+\mu\right)$; setting $\Gamma=e^{g}$ with $g$ the growth rate of \theory{Section}{The circular balanced growth path} gives the circular balanced growth path, the state~C configuration of the classification, including $p^{\mathcal W}=\Theta_{\mathcal W}\Psi$. The code takes $\Gamma$ as an input and defaults to $e^{g}$, since the classification of \theorynote{} is stated for trending technology alone; $\Gamma=1$ is set only for the benchmark check of Section~\ref{sec:q:sol:verification}.

\emph{It gets the closed loop right without being told about it.} The stockpile row is closed with its \emph{primitive} survival factor $s=1-\mu$ and its primitive dividend, and the substitution $\tau^W=-p^{\mathcal W}$ then does the rest: solving~\eqref{eq:q:terminal} for that row is algebraically identical to closing the net waste-stock recursion of \theory{Section}{Costate equations} with the effective factor $1-\mu\left(1-\alpha\varpi\right)$. As the simulated loop closes, $\alpha\varpi\to1$ and the terminal price approaches the perpetuity value automatically. No switch between ``collapse mode'' and ``circular mode'' is needed anywhere in the code, which matters, because which mode the economy is in is precisely what the model is being asked.

\emph{It degenerates gracefully.} On a path with $N_T=0$ the reserve dividend is zero and~\eqref{eq:q:terminal} returns $p^S_T=0$; with $D_T=0$ it returns $p^X_T=0$, which is the exact result of \theory{Section}{Optimality conditions} rather than an approximation.

The residual error introduced by~\eqref{eq:q:terminal} decays with the horizon at the rate at which the economy converges to its long-run configuration, and Section~\ref{sec:q:sol:horizon} reports how it is checked.

\subsection{Horizon}\label{sec:q:sol:horizon}

The horizon has to be far longer than the period of interest, for a reason specific to this model. The economically interesting window is roughly $1850$--$2200$, but the objects that decide the long-run classification --- the exhaustion of the discovery frontier, the closing of the recycling loop, the drift of the anthropogenic stocks towards $\mathcal M_\infty$ --- operate on the residence times $1/\mu$ and $1/\delta$ and on the cumulative bounds of \theory{Section}{The material budget and four accounting facts}, which bind only in the limit. A terminal condition imposed at $2200$ is therefore imposed exactly where the model's own dynamics are least settled, and it contaminates the window of interest through the costates, which are forward sums.

The baseline reports two diagnostics: the change in the window of interest when $T$ is extended, and the residual of the exact costate recursion at $T-1$ evaluated with the terminal closure at $T$. The first is the economically relevant test and the second localises any failure.

Reaching a long horizon is not free, and how it is done matters. A cold starting path deteriorates as $T$ grows --- the simple rules that generate it drift further from the solution the longer they run --- so the horizon is approached by continuation: solve short, extrapolate the tail from growth rates measured \emph{away} from the terminal date, re-simulate the states so the transitions hold exactly, and re-solve. The qualification about measuring growth away from the terminal date is not a detail: the last periods of a finite-horizon solve are distorted by the terminal closure, investment most of all, and extrapolating their period-to-period ratios produces an explosive tail.

There is also a horizon beyond which the system is genuinely stiff, and it is worth naming because it is a property of the model rather than of the solver. On a path that ends the material era by exhaustion, the reserve is driven towards zero and the extraction cost carries the factor $\left(S_{\mathrm{ref}}/S\right)^{\mu_N}$, which diverges there. The Newton system inherits that conditioning, and continuation slows and then stalls as $S$ approaches zero. The stalling horizon is thus a calibration-dependent quantity --- it is where the economy runs out of ore --- and the implementation reports the longest horizon it reached rather than failing silently. Where the window of interest has already stopped moving by that horizon, as it does on the illustrative calibration, nothing is lost.

\subsection{Numerical method}\label{sec:q:sol:method}

\smalltitle{Stacked Newton.} The residual system of Section~\ref{sec:q:model:residuals} is solved as one large square nonlinear system in all $12\left(T+1\right)+6T$ unknowns, rather than by shooting. Shooting is unusable here: with six states and six costates the two-point boundary value problem is stiff, and the corner switches make the shooting map discontinuous. The stacked system, by contrast, has a Jacobian that is \emph{block-tridiagonal in time} --- period-$t$ residuals involve variables dated $t$ and $t+1$ only, a direct consequence of the acyclic within-period ordering of Section~\ref{sec:q:model:within} --- and is therefore cheap to form and to factor.

\smalltitle{Jacobian by structured differencing.} The block structure is exploited by colouring: perturbing the unknowns of every third period simultaneously never mixes two entries of the same Jacobian row, so the full sparse Jacobian is recovered in a number of residual evaluations proportional to the number of unknowns per period rather than to their total. The resulting sparse matrix is factorised directly. This keeps the implementation free of automatic-differentiation dependencies while costing a fixed multiple of one residual evaluation per Newton step.

\smalltitle{Smoothing the switches.} The min-map~\eqref{eq:q:mcp} is piecewise smooth, and a Newton step that lands exactly on a kink can cycle. The solver therefore replaces $\mathrm{clamp}$ by a smoothed version with a smoothing width $\epsilon$, solves a sequence of problems with $\epsilon$ driven to zero, and finishes with $\epsilon=0$ so that the returned solution satisfies the exact complementarity conditions. Because the smoothed problems differ only in one scalar, each is a good starting point for the next, and the sequence typically converges in two or three Newton steps after the first.

\smalltitle{Continuation.} The natural continuation parameters are the ones that make the problem hard. A path is first solved with $\bar R=0$, a hard ceiling well below one, and no trends, which is a well-behaved neoclassical problem; the floor, the ceiling and the trend rates are then raised to their target values in steps, each solved from the previous solution. Policy counterfactuals continue in the dials $\left(\phi^W,\phi^z,\phi^P,\phi^X\right)$ from the planner corner, which is both the best-conditioned point of the policy space and, by Remark~\ref{rem:q:nesting}, the one whose solution is independently characterised.

\subsection{Paths that hit the floor}\label{sec:q:sol:collapse}

A path on which $R_t$ falls below $\bar R$ is not a numerical failure but the model's central prediction in the collapse cells, and it needs handling that does not quietly rule it out. The technology is undefined below the floor, so the smooth system of Section~\ref{sec:q:model:residuals} cannot be continued through it; what replaces it is the discrete comparison of \theory{Section}{The shutdown and its aftermath}. The solver treats the shutdown date $T^{\dagger}$ as a discrete parameter: for each candidate it solves the smooth system on $\left[0,T^{\dagger}\right]$ with the terminal value function $V^{\mathrm{stop}}\left(K,\mathcal W,P\right)$ derived there in place of the terminal conditions of Section~\ref{sec:q:sol:terminal}, and then maximises over $T^{\dagger}$. Because $V^{\mathrm{stop}}$ is available in closed form up to a one-dimensional recursion in the pollution stock, the search is a scalar grid search over dates and adds no dimension to the problem.

Two features of this treatment are worth marking. First, it makes survival and collapse commensurable: a path that closes the loop and a path that shuts down in $2170$ are evaluated with the same welfare criterion, which is exactly the comparison the taxonomy of \theory{Section}{Taxonomy} cannot make analytically. Second, it is where utility degeneracy would bite if consumption reached zero; the closed-form post-shutdown policy of \theory{Section}{The shutdown and its aftermath} shows it does not, consumption declining geometrically at the factor $\left[\left(1-\delta\right)/\left(1+\rho\right)\right]^{1/\eta}$, and $V^{\mathrm{stop}}$ is finite exactly under the collapse ranking condition $\left(1-\delta\right)^{1-\eta}<1+\rho$, one of the explicit conditions of \theory{Section}{Scope: maintained long-run assumptions}. The condition holds automatically for $\eta\le1$ and bounds $\rho$ from below for $\eta>1$; the calibration checks it.

\subsection{Verification}\label{sec:q:sol:verification}

Six checks are run on every solved path, and they are reported rather than merely passed.
\begin{enumerate}
\item \emph{Mass balance.} The collected ledger of \theory{Section}{Aggregate material balances} holds period by period to machine precision, and its cumulated form reconciles the retained endowment $\mathcal M_\infty$ of \theory{Section}{Perpetual circular growth (C)} with the sum of inflows less leakage. This is an identity in the code, so a violation indicates a coding error rather than a solver failure --- which is precisely why it is worth checking.
\item \emph{Planner--market equivalence.} The residual vectors of the two formulations agree to machine precision at randomly drawn points of the state and control space when the dials are at $\left(1,1,1,1\right)$, which tests the implementation proposition of \theorynote{} as an algebraic identity and not merely as a property of the solution.
\item \emph{Analytic long-run blocks.} Two closed forms are compared with the solver, one outside the classification and one inside it. The first is a benchmark test: with $g_A=g_B=0$, the choke case $\varsigma=\infty$ and the terminal closure at $\Gamma=1$, a path started at, and shocked away from, the closed-form dematerialized rest point of the stationary benchmark, \theory{Section}{The dematerialized rest point: the choke case $\varsigma=\infty$}, returns to it. The second is the classification's own check: with trending technology and a soft ceiling, the path converges to the normalised price block and the turnover law of \theory{Section}{The circular balanced growth path}, both solved independently and compared.
\item \emph{Cumulative bounds.} Realised cumulative extraction respects the cumulative extraction bound of \theory{Section}{Optimality conditions}, realised cumulative leakage respects the leakage lemma, and under a hard ceiling realised cumulative throughput respects the throughput lemma, both in \theory{Section}{The material budget and four accounting facts}. These are inequalities the model should satisfy without being told to, and they are a strong test of the accounting.
\item \emph{Horizon and terminal sensitivity}, as described in Section~\ref{sec:q:sol:horizon}.
\item \emph{Local optimality.} The solver returns a point satisfying the necessary conditions, and \theory{Appendix}{Sufficiency} establishes when those are sufficient. Where the conditions of that appendix are not met --- a finite extraction choke, multiplicative damages, a positive $\phi^I$ --- the code runs the feasible-direction test of \theory{Section}{Local sufficiency and numerical verification}: it perturbs the controls at arbitrary dates, lets consumption absorb the goods constraint, and checks that no deviation raises the welfare criterion stated there. A deviation that does is reported as an explicit alternative path rather than as a failure, since in a non-convex problem that is the informative output.
\end{enumerate}
